\section{ChatGPT Agent: un puente entre la investigación y la práctica}

Kimi K2 muestra el pipeline de entrenamiento de un modelo; esta sección pasa a los sistemas de producto, y la pregunta central es cómo un Agent ya entrenado entra en el flujo de trabajo del usuario y asume los riesgos de operar en entornos reales. El segundo material es el anuncio de lanzamiento de ChatGPT Agent. OpenAI integra Operator, Deep Research y las capacidades de conversación de ChatGPT en un sistema de agentes unificado. Este sistema completa tareas mediante su propia computadora virtual, un navegador de texto, un navegador visual, una terminal y acceso a API, y al mismo tiempo conserva la capacidad del usuario de confirmar, tomar el control e interrumpir.

\begin{figure}[H]
\centering
\includegraphics[width=0.96\textwidth]{chatgpt-agent-unified-system.png}
\caption{Sistema unificado de ChatGPT Agent: Operator + Deep Research + ChatGPT + computadora con herramientas. Diagrama conceptual de elaboración propia, basado en la página oficial de OpenAI y en la conversación de 00:43:50--01:53:38.}
\end{figure}

\begin{knowledgebox}{Lectura de la figura: el sistema de producto tiene muchas más capas que el modelo}
ChatGPT Agent no es solo un modelo: incluye navegación, búsqueda, ejecución de código, manejo de archivos, conectores, confirmación del usuario y políticas de seguridad. Estas capas determinan si puede manejar flujos de trabajo reales.
\end{knowledgebox}

\begin{knowledgebox}{Pila de herramientas de ChatGPT Agent}
\begin{tabularx}{\textwidth}{p{0.22\textwidth}p{0.34\textwidth}X}
\toprule
Herramienta/capa & Función & Control de riesgos \\
\midrule
Navegador visual & Interactúa con páginas web hechas para personas: hace clic, filtra, llena formularios & El usuario puede tomar el control del navegador. \\
Navegador de texto & Lee con eficiencia grandes cantidades de texto web & Requiere filtrado de fuentes y citas. \\
Terminal & Ejecuta código y procesa archivos y datos & Las operaciones de alto riesgo deben restringirse. \\
Connectors/API & Se conecta a Gmail, GitHub, calendario, etc. & Control de permisos y de privacidad. \\
User confirmation & Confirma antes de ejecutar operaciones importantes & Reduce las consecuencias de una operación errónea. \\
\bottomrule
\end{tabularx}
\end{knowledgebox}

\subsection{Seguridad: nuevas capacidades traen nuevos riesgos}

La página de OpenAI enfatiza que este nuevo tipo de Agent puede ejecutar acciones en la web, por lo que trae problemas de inyección de prompts, operaciones erróneas, privacidad y tareas de alto riesgo. Las medidas de control incluyen confirmación de operaciones importantes, modo de toma de control, modo de supervisión, rechazo de tareas de alto riesgo, restricción del acceso a datos y toma de control segura del navegador.

\begin{warningbox}{El cambio central en la seguridad de los Agents}
Cuando un modelo de chat dice algo incorrecto, las consecuencias suelen quedarse en el nivel del texto; cuando un Agent ejecuta una acción incorrecta, puede afectar cuentas, archivos, correos, compras, código y operaciones reales del negocio. Por ello, la seguridad se amplía de la moderación de contenido a los permisos, las herramientas, la confirmación y la auditoría.
\end{warningbox}

\begin{knowledgebox}{Niveles de riesgo de seguridad}
\begin{tabularx}{\textwidth}{p{0.20\textwidth}p{0.34\textwidth}X}
\toprule
Riesgo & Ejemplo & Enfoque de protección \\
\midrule
Inyección de prompts & Instrucciones ocultas en una página web inducen al Agent a filtrar datos & Detección, aislamiento, confirmación de acciones importantes. \\
Operación errónea & Enviar un correo, comprar o borrar un archivo por error & Confirmación humana, toma de control, niveles de permisos. \\
Fuga de privacidad & Se leen datos de conectores o de sitios con sesión iniciada & Mínimo privilegio, limpieza de sesiones, desactivación de conectores no relacionados. \\
Tareas de alto riesgo & Operaciones financieras, biológicas o peligrosas & Políticas de rechazo, modo de supervisión, protecciones específicas. \\
\bottomrule
\end{tabularx}
\end{knowledgebox}

\subsection{Resumen de la sección}

Lo central de ChatGPT Agent es un sistema de producto unificado: reúne en un mismo flujo de trabajo las capacidades de investigación, la operación en la web, la ejecución de herramientas y el control del usuario, y muestra el camino de los Agents desde el experimento hasta el producto.
